\begin{center}
{\Large\bfseries Машинное обучение для выбора численного метода\par}
{\large решения линейных дифференциальных уравнений второго порядка\par}
\vspace{3pt}
Третья часть научно-исследовательского отчёта по проекту №2488
\end{center}
\vspace{-4pt}
\section{Исследовательская задача}
Численное решение дифференциального уравнения определяется не только выбранной схемой, но и свойствами коэффициентов, правой части, решения и значениями параметров дискретизации. Одинаковое число итераций может оказаться достаточным для одного уравнения и недостаточным для другого. Стохастический метод дополнительно зависит от случайного набора траекторий. Поэтому фиксированный выбор метода и вычислительного бюджета не обязан быть эффективным для всего рассматриваемого класса задач.

В предыдущих частях исследованы метод Ремизова и метод Монте-Карло для одной резольвентной постановки. Настоящая работа посвящена возможности предварительно выбирать способ вычисления по характеристикам заданного уравнения. Требуется учитывать одновременно точность и время: быстрый, но неточный ответ не решает поставленную задачу, а заведомо избыточные вычисления повышают стоимость точной доставки. При этом существование единственного лучшего метода для каждого уравнения не предполагается: возможны несколько допустимых вариантов, отсутствие точного ответа среди испытанных приближений и неопределённость стохастического результата.

Цель исследования --- построить и независимо проверить обучаемое правило выбора между двумя приближениями Ремизова и одним приближением Монте-Карло. Для этого сформированы признаки коэффициентов, обучены прогнозы достижения допуска и времени, исследована калибровка вероятностей, выбрана политика по семейной перекрёстной проверке и проведено сопоставление с постоянными стратегиями. Научный результат оценивается по фактическим точным доставкам и затратам, а не по обязательному превосходству обучаемого правила.

\section{Математическая постановка и критерий точности}
Рассматривается исходное уравнение
\begin{equation}\label{eq:original}
a(x)y''(x)+b(x)y'(x)+c(x)y(x)=g(x),\qquad x\in\R.
\end{equation}
Здесь $x$ --- независимая переменная, $y$ --- искомая функция; $a,b,c$ --- коэффициенты, $g$ --- заданная правая часть. Для исследуемых гладких ограниченных функций имеются глобальные оценки
\begin{equation}\label{eq:conditions}
0<\underline a\le a(x)\le\overline a,\quad |b(x)|\le B,
\quad |c_R(x)|\le C,\quad \lambda>C,
\end{equation}
где $\underline a,\overline a,B,C$ --- известные числовые границы, $\lambda$ --- параметр резольвенты. Границы не обязаны быть точными экстремумами. Связь с исходным уравнением задаётся равенствами
\begin{equation}\label{eq:operator}
c_R=c+\lambda,\quad g_R=-g,\quad
L=a(x)\partial_x^2+b(x)\partial_x+c_R(x),\quad (\lambda-L)y=g_R.
\end{equation}
Последнее равенство после умножения на $-1$ возвращает \eqref{eq:original}. Ищется решение, ограниченное на всей прямой. Это дополнительное условие выбора решения, отличное от начальных условий задачи Коши. Если в данных заданы $y(-1)$ и $y'(-1)$, они получены из известного точного решения и совместимы с ним; произвольные другие начальные условия определяли бы другую задачу.

\clearpage
Для оператора $L$ используется резольвентное представление
\begin{equation}\label{eq:resolvent}
y(x)=\int_0^\infty e^{-\lambda t}(e^{tL}g_R)(x)\,dt,
\qquad y_{n,T}(x)=\int_0^T e^{-\lambda t}(S(t/n)^ng_R)(x)\,dt,
\end{equation}
где $e^{tL}$ --- полугруппа, $S$ --- аппроксимирующий оператор Чернова, $t$ --- переменная интегрирования, $T$ --- предел усечения, $n$ --- число последовательных применений $S$. В оператор входит именно $c_R$, а не $c$. При Монте-Карло то же представление оценивается по случайным траекториям; диффузия имеет вид $dX_t=b(X_t)dt+\sqrt{2a(X_t)}\,dW_t$, где $W_t$ --- стандартное броуновское движение. Подробные схемы исследованы ранее; здесь они не меняются.

Сравнение выполняется на отрезке $[-1,1]$ в фиксированных узлах $x_i=-1+i/16$, $i=0,\ldots,32$. Для сохранённого приближения $\widetilde y_j$ определяются
\begin{equation}\label{eq:error}
E_j=\max_{0\le i\le32}|\widetilde y_j(x_i)-y_{\mathrm{exact}}(x_i)|,
\qquad Z_j=\ind\{E_j\le\varepsilon\},\qquad \varepsilon=0{,}01.
\end{equation}
Индекс $j$ обозначает метод вместе с его фиксированными параметрами, $Z_j$ --- индикатор точной доставки. В материалах нет отдельного физического обоснования величины $0{,}01$; это установленный критерий эксперимента. Он относится к максимуму в 33 узлах, не к непрерывной норме на всём отрезке и тем более не на $\R$. Успешное завершение вычисления также не означает $Z_j=1$.

Изучаются Ремизов с $n=128$, Ремизов с $n=256$ и Монте-Карло с $M=4096$ траекториями на каждую пространственную точку. При Ремизове использованы 513 узлов квадратуры, $\Delta x=0{,}0125$, кубические сплайны и нулевое продолжение вне рабочего отрезка $[-13,13]$. Этот отрезок не заменяет область исходной задачи. Для Монте-Карло требуемый шаг равен $1/128$, фактический --- $T/N\le1/128$, где $N$ --- число шагов траектории. Удвоение $n$ исследовалось как более затратное уточнение, но не как теорема о двукратном уменьшении ошибки.

\section{Формирование уравнений и экспериментальных данных}
Уравнения строились с заранее известным решением. Сначала задавались гладкие ограниченные композиции $A,B_0,C_0,H$ из констант, синусов, косинусов, гиперболических тангенсов, гауссовых и рациональных функций, их линейных комбинаций и произведений. Затем применялись преобразования
\begin{equation}\label{eq:construction}
a=A_m+A_h\tanh A,\quad b=B\tanh B_0,\quad
c_R=C\tanh C_0,\quad \lambda=C+\delta,\quad\delta>0.
\end{equation}
При $A_m>A_h>0$ они обеспечивают \eqref{eq:conditions} на всей прямой. Различались четыре уровня сложности с растущими глубиной, числом узлов и диапазоном частот композиционных деревьев и пять режимов коэффициентов: постоянный, медленно меняющийся, умеренный, осциллирующий и контрастный. Это способы конструирования разнообразных задач, а не доказанная шкала трудности для каждого метода.

Точное решение имело компактный носитель:
\begin{equation}\label{eq:manufactured}
y_{\mathrm{exact}}(x)=sW_{\mu,\rho}(x)H(x),\quad
W_{\mu,\rho}(x)=\begin{cases}\exp\!\left(1-\dfrac1{1-((x-\mu)/\rho)^2}\right),&|x-\mu|<\rho,\\0,&|x-\mu|\ge\rho.\end{cases}
\end{equation}
Здесь $\mu$ --- центр носителя, $\rho>0$ --- его радиус, $s>0$ --- масштаб. Правая часть получалась аналитической подстановкой $g=ay''_{\mathrm{exact}}+by'_{\mathrm{exact}}+cy_{\mathrm{exact}}$. Так обеспечивалось известное ограниченное решение и возможность измерять полную сеточную ошибку. Это изготовление проверочных задач, а не утверждение о возможности аналитически решать произвольное входное ОДУ.

\clearpage
Покажем реальное уравнение с постоянными коэффициентами (пример 1 приложения~А). Обозначим через $a_1,b_1,c_1$ числа
\begin{equation}\label{eq:realconstant}
\begin{aligned}
a_1&=1.983925129927827+0.4948927992340211\tanh(-0.9445099704819986),\\
b_1&=1.2736839759694616\tanh(-1.4320384163375453),\\
c_1&=1.2336798048825266\tanh(-1.0418788348584225)-1.6551040450654098.
\end{aligned}
\end{equation}
Исходная задача имеет вид $a_1y''+b_1y'+c_1y=g_1(x)$. Её точное решение задаётся \eqref{eq:manufactured}, где
\begin{equation}\label{eq:realexact}
\begin{aligned}
s&=0.4496376326792234,\quad \mu=0.12552944688615264,\quad \rho=2.7787272448026363,\\
H(x)&=1.4204653788060875e^{-1.6583964568669665(x+0.7759412582493684)^2}\\
&\quad+1.4499935751532103\cdot0.774948700390804.
\end{aligned}
\end{equation}
Полная замкнутая правая часть $g_1$, её вспомогательные функции и совместимые начальные данные даны в п.~А.1; это то же уравнение, а не его упрощённый аналог. При $n=128$ ошибка равна $0{,}00291937$, при $n=256$ --- $0{,}00146141$: оба приближения допустимы. Монте-Карло в сохранённой реализации дал $0{,}01026241$, то есть нарушил порог, несмотря на близость к нему. Графики Б.1--Б.4 позволяют отделить постоянство коэффициентов от непостоянной правой части и асимметричного решения.

Семейство определялось топологией четырёх композиционных деревьев; изменение числовых параметров внутри одной топологии не создавало независимое семейство. Поэтому случайное разделение отдельных уравнений могло бы перенести родственные конструкции и в обучение, и в проверку. До обучения разделялись семейства; повторные записи одного расчёта удалялись, а независимые реализации Монте-Карло сохранялись как повторные наблюдения одного уравнения.

Обучающие данные точности для $n=128$ содержали 103 случая из 90 семейств, для $n=256$ --- 40 случаев из 40 семейств. Для Монте-Карло имелось 40 уравнений и 180 различных случайных реализаций, из которых 12 достигли допуска. Нельзя считать эти 180 реализаций 180 независимыми уравнениями. Для времени использованы только сопоставимые измерения: 9 для $n=128$, 36 для $n=256$ и 40 для Монте-Карло после агрегации повторов по уравнению и профилю. Отсутствующее время не заменялось вымышленным наблюдением.

\section{Вероятностная модель и правило выбора}
Входом модели служили только заданные коэффициенты $a,b,c,g$ на 256 узлах и их допустимые аналитические характеристики, но не точное решение, его производные, численные ответы или ошибка. Получены 253 признака: 128 выборочных значений функций, 56 статистик уровня, вариации, конечных разностей и спектрального состава, 6 межканальных корреляций, 48 характеристик структуры коэффициентов и 15 глобальных оценок и производных от них масштабов. Каналы $c_R$ и $g_R$ получаются из исходных коэффициентов по \eqref{eq:operator}, поэтому не добавляют скрытой информации о целевом ответе. Структурные признаки описывают только коэффициенты, не дерево известного решения.

Вариация, производные и спектральный состав потенциально характеризуют изменение функций; глобальные границы задают масштаб дрейфа, диффузии и затухания. Однако их информативность для конкретной ошибки исследуется статистически. Само наличие частых колебаний не доказывает большого значения $E_j$.


Обучение и калибровка проверялись трёхблочной перекрёстной проверкой по семействам. Для каждого внешнего блока все его семейства исключались из обучения, внутренней калибровки и выбора политики. Калибратор обучался на отдельных внутренних трёхблочных прогнозах. Среди 20 сохранённых правил сравнивались пороговые варианты, близость вероятностей, стоимость успеха и постоянный выбор. Внутренний эмпирический отбор требовал отсутствия индивидуальных потерь точных доставок относительно постоянного $n=256$, затем максимизировал число точных ответов и минимизировал прогнозируемую стоимость. Это обучающий критерий, не доказательство безопасности на популяции. Окончательная политика выбрана по всем разрешённым обучающим данным; её оценка после такого отбора не выдавалась за независимую. Ранее просмотренные проверочные данные для настройки не использовались.

\clearpage
Для вектора признаков $z$ обучались отдельные прогнозы вероятности успеха
\begin{equation}\label{eq:probability}
p_j(z)=\Pr\{E_j\le0{,}01\mid z,j\}.
\end{equation}
У Монте-Карло вероятность относится также к новой случайной реализации при фиксированном профиле. У детерминированного метода она выражает неопределённость прогноза по классу уравнений, а не случайность повторного исполнения одной задачи. Для обучения использованы ансамбли градиентного бустинга над деревьями (CatBoost): 120 деревьев глубины 2, шаг обучения $0{,}05$ и регуляризация листьев 10. Эти настройки фиксировались, их преимущество перед всеми альтернативными алгоритмами не установлено.

Если из $m_{ij}$ реальных повторов уравнения $i$ точны $k_{ij}$, целевое значение равно $q_{ij}=k_{ij}/m_{ij}$. У Ремизова повторные детерминированные ответы дают одну метку; у Монте-Карло число различных реализаций входит в биномиальную функцию потерь:
\begin{equation}\label{eq:loss}
\mathcal L_j=-\sum_i w_i m_{ij}\{q_{ij}\log p_j(z_i)+(1-q_{ij})\log[1-p_j(z_i)]\}.
\end{equation}
Вес $w_i$ обратен числу уравнений соответствующего семейства среди обучающих строк. Это ограничивает перевес многократно представленных родственных конструкций; множитель $m_{ij}$ учитывает реальное число стохастических наблюдений. Признаки стандартизировались по статистикам только обучающей части каждого разбиения.

Вероятности калибровались регуляризованным сигмоидальным преобразованием
\begin{equation}\label{eq:calibration}
\widehat p_j=\sigma\!\left(\alpha_j\operatorname{logit}(p_j^{\rm raw})+\beta_j\right),\quad
\sigma(v)=(1+e^{-v})^{-1},\quad\operatorname{logit}(p)=\log\frac p{1-p}.
\end{equation}
Перед logit значения ограничивались интервалом $[10^{-6},1-10^{-6}]$. Параметры $\alpha_j,\beta_j$ подбирались по прогнозам на исключённых целых семействах с добавкой $(\alpha_j-1)^2+\beta_j^2$ к потере при $0\le\alpha_j\le10$, $|\beta_j|\le10$. Калибровка предназначена для согласования прогнозов с наблюдениями на доступных данных, но не создаёт сертификата надёжности каждого уравнения.

Время предсказывалось отдельно: общая регрессия для Ремизова использовала цель $\log(t/n)$, а для Монте-Карло --- $\log t$. Поэтому
\begin{equation}\label{eq:time}
\widehat t_{R,n}(z)=n\exp F_R(z),\qquad\widehat t_M(z)=\exp F_M(z).
\end{equation}
Линейное масштабирование по $n$ является предположением модели времени, а не теоремой о точности или стоимости любого алгоритма. Модели не обучались на недостающих вымышленных временах.

Пусть $J=\{R_{128},R_{256},M\}$. Замороженная политика использует
\begin{equation}\label{eq:eligible}
p_* =\max_{j\in J}\widehat p_j,\qquad
\mathcal A=\{j\in J:\widehat p_j\ge p_*-0{,}02\},\quad
\mathcal A_R=\mathcal A\cap\{R_{128},R_{256}\}.
\end{equation}
Если $\mathcal A_R\ne\varnothing$, выбирается $r_*\in\mathcal A_R$ с наименьшим $\widehat t$; при равенстве предпочтён $n=128$. Итоговое правило в точности имеет вид
\begin{equation}\label{eq:policy}
\widehat j=\begin{cases}M,&\mathcal A_R=\varnothing,\\
M,&M\in\mathcal A,\ \widehat t_M\le0{,}9\widehat t_{r_*},\\
r_*,&\text{иначе}.\end{cases}
\end{equation}
Так сохраняются близкие по вероятности варианты и требуется не менее 10\% прогнозной экономии для выбора Монте-Карло вместо подходящего приближения Ремизова. Абсолютного порога $0{,}95$ в этом правиле нет: непустое $\mathcal A$ возможно и при низких вероятностях, поэтому фактический ответ может оказаться неточным.

\clearpage
\section{Методика независимой проверки}
После заморозки модели и правила \eqref{eq:policy} отобраны 60 новых семейств, топологически не пересекающихся с прежними задачами: по три на каждую из 20 комбинаций сложности и режима. Отбор был определён до численных исходов, без поиска благоприятных ошибок. На каждом уравнении измерены оба варианта Ремизова; Монте-Карло выполнен на 20 заранее выбранных задачах, по одной на ячейку. Так получены 140 вызовов. Уменьшение объёма стохастической части сохраняет полное сравнение фактических решений модели с двумя базовыми стратегиями, поскольку все замороженные решения выбрали Ремизова: 21 раз $n=128$, 39 раз $n=256$.

Для основной оценки повторно использовалось одно и то же измерение выбранного приближения, а не дополнительное решение. Измерялось полное время тёплого вызова метода в одном окружении; общей загрузки, подготовки среды и компиляции в нём нет. Для модели добавлено ранее измеренное время получения признаков и прогноза. Поэтому оценка не равна всей задержке возможного практического развёртывания. Порядок измерений сбалансирован; один замер на профиль не оценивает его повторный временной разброс.

\section{Результаты независимого эксперимента}
\begin{table}[!ht]\centering\caption{Точность и вычислительная стоимость на 60 новых семействах}\label{tab:main}
\begin{tabular}{lrrr}\toprule Метод выбора&Точных ответов&Время, с&Отказов\\\midrule
Машинное обучение&57 из 60&1164,27&0\\
Ремизов, $n=256$&57 из 60&1411,01&0\\
Ремизов, $n=128$&53 из 60&706,41&0\\\bottomrule\end{tabular}\end{table}
Результат 57 из 60 означает, что в трёх сохранённых ответах максимум \eqref{eq:error} превысил допуск. Отказов и вычислительных сбоев не было; неточные ответы также оплачены. Суммарное время модели включает $0{,}103670$ с ранее измеренных затрат на признаки и прогноз. Относительная экономия к постоянному $n=256$ равна
\begin{equation}\label{eq:saving}
\frac{1411{,}006868253-1164{,}270215714}{1411{,}006868253}\cdot100\%
=17{,}4866\%\approx17{,}49\%.
\end{equation}
\fig{summary_02_cost_accuracy}{fig:cost}{.86}{Число точных доставок и стоимость на одной финальной выборке. Самый дешёвый постоянный вариант менее точен; обучаемый выбор не доминирует обе стратегии одновременно.}

\clearpage
Относительно постоянного $n=128$ модель получила ещё четыре точных ответа, сохранив все его 53 точные доставки, но затратила на $64{,}82\%$ больше времени. Стоимость на точную доставку составила соответственно $20{,}426$ с для модели, $24{,}755$ с для $n=256$ и $13{,}328$ с для $n=128$. Разные знаменатели здесь существенны: дешёвые неточные ответы нельзя исключать из оплаченных затрат.

Распределения ошибок на рис.~\ref{fig:ecdf} показывают, что одинаковое число точных ответов не означает одинаковые ошибки. Выбор более дешёвого допустимого приближения может увеличивать ошибку, не пересекая порог. Монте-Карло анализируется отдельно: все 20 вычислений завершены, но ни одно не достигло допуска. Для остальных 40 его точность неизвестна. Нельзя ни смешивать знаменатели 20 и 60, ни считать неизмеренные варианты неудачными.
\fig{summary_01_error_distributions}{fig:ecdf}{.96}{Эмпирические распределения сеточной ошибки. Для Ремизова и модели использованы 60 задач, для Монте-Карло --- 20. Вертикаль соответствует допуску $0{,}01$.}

Неопределённость оценивалась парным повторным отбором семейств; подробности приведены в приложении~В. Вырожденный интервал разности точности с $n=256$ не доказывает равенства рисков. Наблюдаемые $57/60=95\%$ не гарантируют надёжность следующего ответа $95\%$.

\fig{summary_04_selection_outcomes}{fig:choice}{.90}{Выбор модели относительно двух измеренных приближений Ремизова: 19 экономных точных выборов, 4 необходимых уточнения, 34 избыточно дорогих точных ответа и 3 неточных.}

\clearpage
\section{Характерные уравнения и неоптимальные решения}
Пример 1 с уравнением \eqref{eq:realconstant} и точным решением \eqref{eq:realexact} показывает достаточность меньшего вычислительного бюджета. Коэффициенты постоянны, но правая часть и решение меняются. Сравнение Б.1--Б.4 показывает, что визуальное совпадение кривых недостаточно для принятия решения: отдельная MC-ошибка чуть выше $0{,}01$ остаётся нарушением допуска.

В примере 3, заданном уравнением~\eqref{eq:ex3-ode}, коэффициенты меняются плавно, а решение имеет вид \eqref{eq:manufactured} с $s=6.034151780918269$, $\mu=-0.03666028004499061$, $\rho=1.7373814842542066$. Его функция формы равна
\begin{equation}\label{eq:example3shape}
\begin{aligned}
H_3(x)&=e^{-\cos^2(0.162537284342751x+0.4547833238272104)}\\
&\quad\times\bigl[-0.8980719428395056\tanh(-0.3613725371874747x-1.1967457073726797)\\
&\qquad-0.6907321591909864e^{-0.46929024038993034(x-1.3123761993740666)^2}\bigr].
\end{aligned}
\end{equation}
Исходное уравнение здесь $a_3y''+b_3y'+c_3y=g_3$ с явными числовыми коэффициентами п.~А.3; оно не заменяется уравнением с усреднёнными коэффициентами. Уточнение действительно пересекает допуск, хотя различие кривых решений на общем масштабе мало (рисунки Б.9--Б.12).

Для контраста пример 5 с уравнением~\eqref{eq:ex5-ode} имеет несколько экстремумов и меняющее знак решение $y=sW_{\mu,\rho}H_5$, где $s=0.4266912316222408$, $\mu=0.20899798693774746$, $\rho=1.9721061459030749$; полное числовое определение $H_5$, коэффициентов и правой части приведено там же. Оба приближения нарушают порог. Уточнённая ошибка на части узлов допустима, но её максимум всё ещё велик (Б.17--Б.20).
\begin{table}[!ht]\centering\caption{Три характерных случая; полные постановки --- в приложении А}\label{tab:cases}
\begin{tabular}{clrr}\toprule Пример&Наблюдение&$E_{128}$&$E_{256}$\\\midrule
1&Оба приближения точны&0,00291937&0,00146141\\
3&Уточнение достигает допуска&0,01022150&0,00376125\\
5&Оба приближения неточны&0,04272617&0,02179748\\\bottomrule\end{tabular}\end{table}

В целом на 57 задачах с хотя бы одним точным приближением Ремизова модель выбрала самое быстрое точное приближение лишь в 23 случаях. В 34 случаях она предпочла более дорогой $n=256$ при точном $n=128$. Матрица на рис.~\ref{fig:choice} также выделяет три задачи, где оба варианта неточны. Следовательно, экономия к одной базовой стратегии не равнозначна оптимальному выбору на каждом уравнении. Форма коэффициентов на рисунках описывает наблюдаемый пример, но не устанавливает причину ошибки или порядок сходимости.
\clearpage
\section{Обсуждение и научные выводы}
Независимая проверка подтвердила ограниченный, но содержательный эффект: на 60 новых семействах обучаемое правило сохранило все 57 точных доставок постоянного $n=256$ и уменьшило измеренную стоимость на $17{,}49\%$. Результат не является доминированием обеих базовых стратегий: $n=128$ дешевле, но даёт 53 точных ответа. На рис.~\ref{fig:times} приведены парные измерения и накопленная стоимость без утверждения об устойчивом времени каждого нового запуска.
\fig{summary_03_time_comparison}{fig:times}{.96}{Сопоставимые времена и накопленная стоимость по фиксированному порядку 60 задач. Повторное использование одного измерения сохраняет парность сравнения стратегий.}

Математически и вычислительно реализованы единая резольвентная постановка, построение проверочных уравнений с известным решением, прогнозы вероятности точности и времени, семейная проверка и правило выбора трёх приближений. При этом требование $E_j\le0{,}01$ не обеспечено для каждого ответа: модель и постоянный $n=256$ нарушили его на трёх задачах. Положительный перенос на другие классы ОДУ, непрерывную норму ошибки или произвольные начальные условия данным исследованием не установлен.

Поддержка Ремизова и Монте-Карло присутствует в модели, однако в финальном опыте Монте-Карло не выбран ни разу. Его 20 проверенных реализаций были неточны, но $0/20$ не доказывает принципиальной неэффективности метода: исследованы конкретные число траекторий, шаг и задачи. Исторический пример 7 (п.~А.7, рисунки Б.25--Б.29) содержит один точный и быстрый MC-ответ из четырёх; он иллюстрирует случайную изменчивость, не добавляется к независимым данным и не подтверждает надёжность следующего запуска.

Дальнейшее исследование обоснованно связано с ранжированием точных приближений, отказом при малой прогнозной допустимости и независимым разделением дисперсии и смещения Монте-Карло. После раскрытия финальных данных модель не настраивалась; новые решения не вычислялись.

Таким образом, реализация обучаемого выбора завершена, а её эмпирическое качество относительно двух постоянных стратегий измерено. Более сильная исходная цель --- надёжно выбирать самый быстрый допустимый метод для каждого нового уравнения, включая полезный выбор Монте-Карло, --- полностью научно не подтверждена. Исследовательская версия не объявляется гарантированным средством точной доставки.

\vspace{4pt}\noindent\textbf{Предшествующие части исследования.} Исследование и оптимизация численной реализации метода Ремизова. Отчёт по проекту №2488, 2026, 10~с.; Сравнение численных методов решения линейных ОДУ и прогнозирование их точности с помощью машинного обучения. Отчёт по тому же проекту, 2026, 6~с.
